\documentclass[12pt]{article}
\usepackage{amsmath, amssymb, amsthm}

\title{VGASP Snowman: Invariant Proofs}
\author{Borealis S. Hedling}
\date{Version 1.0 (2026-10-08)}

\begin{document}
\maketitle

\section*{Overview}
This document provides formal proofs for invariant behavior in the VGASP Snowman
system. Invariants are governed by the invariant tensor $I_{ij}$, which encodes
identity, symmetry, and canonical surface constraints. We prove:

\begin{itemize}
    \item invariants are preserved under roll
    \item invariants restrict stack geometry
    \item invariants forbid premature decoration
    \item invariants couple deterministically to curvature and holonomy
\end{itemize}

All proofs rely on tensor definitions in \texttt{vgasp-snowman-math.tex}.

\section{Preliminaries}

\subsection{Invariant Tensor}
The invariant tensor $I_{ij}$ is defined as a rank-2 symmetric tensor satisfying:
\[
I_{ij} = I_{ji}, \qquad \nabla_k I_{ij} = 0.
\]
The covariant constancy condition enforces identity preservation across all
allowed Snowman operations.

\subsection{Operator Definitions}
Let $\mathcal{R}$, $\mathcal{S}$, and $\mathcal{D}$ denote the roll, stack, and
decorate operators respectively. Each operator acts on a geometric state
$\Sigma$.

\[
\mathcal{R}: \Sigma \rightarrow \Sigma', \qquad
\mathcal{S}: \Sigma' \rightarrow \Sigma'', \qquad
\mathcal{D}: \Sigma'' \rightarrow \Sigma'''.
\]

\section{Invariant Preservation Under Roll}

\begin{lemma}
Roll preserves invariant surfaces.
\end{lemma}

\begin{proof}
Roll modifies curvature $K_{ij}$ and holonomy $H_{ijk}$ but does not alter
$I_{ij}$. Since $\nabla_k I_{ij} = 0$, any deformation induced by roll must
satisfy:
\[
\mathcal{L}_{\mathcal{R}} I_{ij} = 0,
\]
where $\mathcal{L}$ denotes the Lie derivative. Because roll is defined to be
tangent to invariant surfaces, the induced flow preserves $I_{ij}$.
\end{proof}

\section{Invariant Restrictions on Stack}

\begin{lemma}
Stacking is only allowed on curvature-invariant surfaces.
\end{lemma}

\begin{proof}
Stack requires alignment of centers of mass. Let $n^i$ be the surface normal.
Curvature-invariant surfaces satisfy:
\[
I_{ij} n^i n^j = \text{constant}.
\]
If stack occurs on a non-invariant surface, the identity constraint is violated:
\[
\delta I_{ij} \neq 0.
\]
Since $\nabla_k I_{ij} = 0$, any non-zero variation contradicts the definition
of invariant surfaces. Therefore stack is forbidden unless the surface satisfies
the curvature-invariant condition.
\end{proof}

\section{Invariant Prohibition of Premature Decoration}

\begin{theorem}
Decoration cannot occur before stack.
\end{theorem}

\begin{proof}
Decoration modifies surface geometry via a local operator $\mathcal{D}$ acting
on $\Sigma''$. Let $\gamma$ be a decoration curve. Decoration requires:
\[
I_{ij} \gamma^i \gamma^j = \text{constant}.
\]
Before stack, curvature $K_{ij}$ and holonomy $H_{ijk}$ are not stabilized.
Thus $\gamma$ is not guaranteed to lie on an invariant surface. If decoration
occurs prematurely, the identity constraint is violated:
\[
\mathcal{L}_{\mathcal{D}} I_{ij} \neq 0.
\]
Therefore decoration is forbidden until stack stabilizes invariant surfaces.
\end{proof}

\section{Invariant Coupling to Curvature and Holonomy}

\begin{theorem}
Invariant surfaces couple deterministically to curvature and holonomy.
\end{theorem}

\begin{proof}
Curvature-invariant surfaces satisfy:
\[
I_{ij} K^{ij} = \kappa,
\]
for constant $\kappa$. Holonomy-invariant surfaces satisfy:
\[
I_{ij} H^{ijk} = 0.
\]
Since both conditions derive from contractions of $I_{ij}$ with other tensors,
and since $I_{ij}$ is covariantly constant, the coupling is deterministic:
\[
\nabla_m (I_{ij} K^{ij}) = 0, \qquad
\nabla_m (I_{ij} H^{ijk}) = 0.
\]
Thus invariant surfaces constrain curvature and holonomy in a manner that is
independent of operator history.
\end{proof}

\section*{Conclusion}
We have shown that invariant surfaces:

\begin{itemize}
    \item are preserved under roll
    \item restrict stack geometry
    \item forbid premature decoration
    \item couple deterministically to curvature and holonomy
\end{itemize}

These results form the foundation for dissipation proofs and JSON operator
validation.

\end{document}
